Before there is a harvest, someone has to plant. Before anyone carries away a basket of grapes, someone has cleared the ground, set the vines, and tended them. Jesus begins today's parable with that work already done. The vineyard has a hedge, a winepress, and a tower. The tenants enter a place prepared for them. Their work matters, but they did not create the gift that makes their work possible.

That is where we enter the story, too. We have received life, people who have cared for us, opportunities to learn, and the Gospel itself. Even when much has been painful or lacking, we are not our own creators. What have we received, and what is growing from it? That question reaches further than a list of accomplishments. It asks what our gifts are becoming in our dealings with other people.

Isaiah makes the question painfully clear. His song begins with a beloved vineyard and the care lavished upon it. But the expected grapes do not come. At the end he tells us what the image means. God's vineyard is Israel, and the fruit he seeks is justice. Instead, the cry of people wronged rises from it. The problem is not that the vineyard looks unimpressive. The problem is what life inside it has become.

We can recognize that danger. A home can be orderly while one person is never heard. A workplace can succeed while someone is treated unfairly. Religious activity can fill a calendar while a neighbor's need goes unanswered. The question is not simply whether something is growing. Resentment grows. Greed grows. The desire to have the last word grows. God asks for the fruit of justice and love.

Jesus' parable brings another part of the danger into view. Here the tenants turn against the owner. They injure his messengers and finally kill his son because they want the inheritance. They want to keep the gift while refusing the giver. Saint John Chrysostom notices both how much the owner has already provided and how patiently he continues to send. Such generosity makes the tenants' response more serious. It gives them every reason to answer differently.

Jesus is confronting leaders who resist him. He is not giving Christians a reason to blame every Jewish person or to congratulate ourselves. The new stewards will also owe fruit. We hear this Gospel as people entrusted with gifts and responsible for what we do with them. The temptation to push aside a correction because it threatens our control is close enough to home.

Yet the tenants' violence does not settle the future. Jesus speaks of the rejected stone becoming the cornerstone. He himself is that foundation. The Son they reject is not defeated by their rejection. Our hope rests upon Christ, crucified and risen. A life can be rebuilt upon him even when its own claims to goodness have collapsed. His mercy opens the possibility of repentance; it does not make the wrong harmless.

Perhaps the demand for fruit leaves us afraid. We know some of the wrong we have done. We also know how often good intentions have come to nothing. Does God want a harvest we cannot produce? Must we make ourselves good enough before we can approach him?

The Gospel acclamation directs us to Christ's choice of his disciples and his purpose that they bear lasting fruit. Saint Augustine attends to that order. Christ does not choose them because they have already made themselves fruitful. He chooses them so that they may become fruitful. His mercy comes first. The branch receives life from the vine while it bears its fruit; it never reaches a stage where it can do without that life.

The Collect, the opening prayer of this Mass, places us before God's generosity, which exceeds what we deserve and even what we know how to ask. We approach him as people needing mercy. We do not present a successful harvest as the price of his attention. The demand for a changed life and the confidence to ask for help belong together. We can admit our failure because his generosity is greater than our failure. The entrance antiphon likewise recalls a threatened people's trust in their Creator. We depend upon a power greater than the troubles we cannot master.

Augustine also asks what the fruit is. Jesus goes on to command his disciples to love one another. That is the fruit: charity. It takes shape in patience, kindness, fidelity, and the desire for another person's good. Beside Isaiah's vineyard, that love must include justice. We cannot call ourselves loving while refusing to face the harm we cause. Receiving mercy means becoming willing to let mercy change us.

Saint Paul gives that change an ordinary, practicable shape. Bring your needs to God with thanksgiving. Give your attention to what is true and just. Put into practice what you have received. Thanksgiving recalls that we have already been given something; petition admits that we still need help. Neither requires us to pretend that suffering has disappeared. God's peace is not a test we fail whenever we feel troubled.

Paul's words also reach what we allow to occupy our minds. If we keep rehearsing an insult, we prepare another angry response. If we look patiently for what is true, we may discover a correction we need to receive. Thought prepares action. Paul does not stop with thinking about good things. He asks for deeds. The fruit of a gift becomes visible when gratitude changes what we actually do.

This week, begin with one relationship in which something has been entrusted to you. Consider one person whose good depends, in some way, upon your response. Perhaps you owe an apology without an excuse attached. Perhaps you need to stop repeating an accusation you have not checked. Perhaps someone needs your patient attention or a wrong needs to be put right. Choose one truthful act of love and carry it out. Peace grows through truth and justice; it cannot be built by asking the injured to hide their pain.

And bring that intention with you into the Eucharist. The Lord who calls for fruit gives himself to nourish our life. Our offering looks toward God's sanctifying action: obedience still depends upon his gift. Saint Thomas Aquinas explains that the sacrament's power comes from Christ and his Passion. It sustains spiritual life and stirs charity into action. Communion therefore reaches beyond the moment of receiving. The gift is meant to shape how we live, including how we treat the person we would otherwise avoid.

The Prayer after Communion asks for that transformation. What we receive is to form the receiver. Christ's gift is not an ornament added to a life that remains closed to him. His love reaches our words, our choices, and the use of our strength. The hands that have harmed can begin to serve. The mouth that has defended every fault can tell the truth. These are small beginnings, but they are real beginnings in the life he gives.

Some fruit takes patience. One honest apology does not repair every wound. One act of justice does not undo every injustice. The psalm gives a wounded people words with which to ask God for life. They do not have to restore themselves before they can turn to the one who planted them. We, too, can return to him with both our need and our willingness to change. One of this Sunday's Communion antiphons speaks of seeking and waiting for the good Lord. The alternative speaks of the many sharing one body. We need that patient hope, and we need one another. The life Christ nourishes bears fruit in a love that keeps serving while healing is still incomplete.

Before there is a harvest, someone has to plant. God has not merely stood at a distance waiting for us to become useful. He has given us his Son. The rejected Son is our cornerstone; the living Christ is our nourishment. Receive his correction with hope. Answer his generosity with a concrete act of love. The fruit he seeks grows from the life he gives, and the charity he brings to maturity will outlast the harvests we can count.
